Given the current context in the field of cybersecurity, this Bachelor’s Thesis develops an intrusion detection system based on machine learning, more specifically on the Random Forest algorithm.

Using the CIC-IDS2017 dataset, the network traffic contained in it is processed and a feature engineering process is carried out in order to complement the original data. The resulting dataset is then cleaned and used to train and evaluate the Random Forest classifier, with the aim of classifying each sample either as benign traffic or as a specific type of attack, thus formulating the task as a multiclass classification problem. Finally, the relevance of the defined features and the performance of the model are analysed by proposing two experiments and comparing the results obtained in each of them.

\keywords{intrusion detection systems, machine learning, Random Forest, CIC-IDS2017, feature engineering}
